\section{Root exposure and conditional comparability}
\label{sec:exposure}

Assume for now that $n\ge2k$, so that $u=n-k\ge k$.  Fix $i\le L$ and
condition on an arbitrary value of $C_i$.  Write $R$ for its set of $k$
variables, called the \emph{roots}, and let $w\in\{0,1\}^R$ be the unique
root assignment falsifying $C_i$.

For every $j\le L$, $j\ne i$, first reveal the exact subset of roots used
by $C_j$ and all literal signs on that subset.  Classify the index as
follows.
\begin{itemize}
\item If $C_j$ meets no root, call it \emph{ordinary} and leave all of its
      variables and signs hidden.
\item If $C_j$ meets exactly one root and its root literal is true under
      $w$, call it a \emph{test clause} for that root and leave its $k-1$
      nonroot literals hidden.
\item In every other case, reveal the entire clause.
\end{itemize}
Let $\cE$ be the sigma-field generated by this exposure, including the
clause-type schedule at future indices.  Define
\begin{align}
    Z&=\#\{j\le L:j\ne i,\ C_j\text{ meets }R\},\label{eq:Z-definition}\\
    V&=\#\{j\le L:j\ne i,\ C_j\text{ meets at least two roots}\},
      \label{eq:V-definition}\\
    d&=\max\{1,Z\}.\label{eq:d-definition}
\end{align}
Clauses counted by $V$ are called \emph{collisions}.

Because the clause coordinates are independent and the exposure conditions
separately on each coordinate, the hidden parts remain conditionally
independent across clause indices.  Under the local criterion, their density
bounds are precisely assumptions~(i)--(ii).  The following lemma verifies
those assumptions from the simpler pointwise hypothesis used in
Theorem~\ref{thm:main}.

\begin{lemma}[Pointwise comparison implies conditional comparison]
\label{lem:conditional-comparability}
Under \eqref{eq:uniform-comparability}, the conditional law of each hidden
ordinary clause has density between $a/b$ and $b/a$ relative to $U_{u,k}$.
The conditional law of each hidden test part has density between $a/b$ and
$b/a$ relative to $U_{u,k-1}$.
\end{lemma}

\begin{proof}
For a fixed index $j$, write
\[
    f_j=\frac{\dd\mu_{n,j}}{\dd U_{n,k}}.
\]
The pointwise assumption \eqref{eq:uniform-comparability} gives
$a\le f_j\le b$.  Let $E_j$ denote the root subset and root signs revealed
at index $j$.  Relative to the uniform conditional law of the hidden
nonroot part, the conditional density under $\mu_{n,j}$ is
\[
    \frac{f_j(C_j)}{
      \E_{U_{n,k}}[f_j(C_j)\mid E_j]}.
\]
Both numerator and denominator lie in $[a,b]$, so this ratio lies in
$[a/b,b/a]$.  Under the uniform law, conditioning on no roots leaves a
uniform proper $k$-clause on the $u$ nonroots, while conditioning on a
specified single root and sign leaves a uniform proper $(k-1)$-clause on
the nonroots.
\end{proof}

For $i\le m\le L$, let $S_m\subseteq\{0,1\}^u$ be the set of nonroot
assignments $z$ such that $(w,z)$ satisfies $F_m^{(-i)}$.  The sets $S_m$
decrease with $m$.  Define the filtrations
\begin{align*}
    \cG_m
    &=\cE\vee
      \sigma(\text{hidden ordinary-clause contents through index }m),\\
    \cH_m
    &=\cG_m\vee
      \sigma(\text{hidden test parts through index }m).
\end{align*}
The set $S_m$ is $\cG_m$-measurable because every test clause is already
satisfied by its root literal under $w$, regardless of the hidden nonroot
part.

For $i\le m<L$, put
\begin{equation}
    A_m=
    \{F_m^{(-i)}\text{ is satisfiable and }
      F_m^{(-i)}\wedge C_i\text{ is unsatisfiable}\}.
    \label{eq:A-definition}
\end{equation}
Then
\begin{equation}
    A_m=\{T\le m<T^{(-i)}\},
    \qquad
    D_i=\sum_{m=i}^{L-1}\ind_{A_m}.
    \label{eq:delay-count}
\end{equation}
The event $A_m$ is $\cH_m$-measurable.  On $A_m$, the set $S_m$ is
nonempty, and every solution of $F_m^{(-i)}$ has root assignment $w$.

Finally, put
\begin{equation}
    p_m=q_d(S_m).
    \label{eq:p-definition}
\end{equation}
Then $p_m$ is $\cG_m$-measurable and nondecreasing in $m$.
